\section{Проектиране на системата}
\label{sec:design}

Разделът излага проектните решения на системата --- слоестата структура, договорите,
табличното преобразуване, алгоритмите за разположение, системата от свойства --- заедно с
алтернативата, спрямо която всяко е взето: интерфейсът остава преносим, а платформено
обусловената част е сведена до необходимия минимум.

\subsection{Синтез на общата структура}
\label{subsec:structure}

Системата е разделена на девет слоя, подредени по посока на зависимостите: по-голям номер
може да зависи от по-малък, никога обратното. Правилото не е стилово --- то прави
преносимата част преносима: платформена частност не може да проникне в алгоритъм, който
трябва да дава един и същ резултат навсякъде.

\begin{enumerate}[topsep=2pt,itemsep=0pt,leftmargin=*,label=\arabic*.]
\setcounter{enumi}{-1}
\item \textbf{Графични примитиви} --- цвят, точка, размер, правоъгълник, геометрия, платно.
\item \textbf{Договори} --- интерфейсите на виртуалните изгледи и техните обработчици.
\item \textbf{Обработчици (handler-и)} --- таблиците за преобразуване.
\item \textbf{Платформени слоеве} --- свързването с native инструментариума.
\item \textbf{Разположение} --- алгоритмите за измерване и подреждане.
\item \textbf{Контроли} --- програмен интерфейс, свойства, свързване на данни.
\item \textbf{Декларативен слой} --- преобразуване на описание в обектен граф.
\item \textbf{Услуги на устройството} --- независими от графичния стек.
\item \textbf{Инициализация} --- изграждане на приложението и регистрация.
\end{enumerate}

Структурата е показана на \figref{fig:layers}.

\begin{table}[H]
\caption{Съответствие директория $\rightarrow$ слой $\rightarrow$ обем}
\label{tab:dirs}
\centering\small
\begin{tabular}{@{}L{3.4cm}rL{7.4cm}@{}}
\hline
\textbf{Директория} & \textbf{ЕП} & \textbf{Слой} \\
\hline
\code{graphics/}   & 27  & 0 --- примитиви \\
\code{core/}       & 64  & 1 + 2 --- договори \emph{и} обработчици \\
\code{layouts/}    & 8   & 4 --- алгоритми за разположение \\
\code{controls/}   & 159 & 5 --- програмен интерфейс, свойства, свързване \\
\code{xaml/}       & 35  & 6 --- декларативен слой \\
\code{essentials/} & 40  & 7 --- услуги на устройството \\
\code{hosting/}    & 9   & 8 --- инициализация \\
\code{animations/} & 6   & 5 --- анимации \\
\code{platform/}   & 318 & 3 --- пет платформени слоя \\
\hline
\end{tabular}
\end{table}
\noindent\small ЕП = единици за превод (\code{.cpp}/\code{.mm}); публични заглавни
файлове --- 716.\normalsize

Две наблюдения: обработчиците живеят заедно с договорите, не в отделна директория, защото
се изменят синхронно; и платформените слоеве са най-обемната част --- количествен израз на
твърдението, че платформената обвързаност е неотстранима, не подлежаща на абстрахиране.

\begin{figure}[H]
\centering
\begin{tikzpicture}[font=\footnotesize, node distance=0.22cm,
  lay/.style={draw, rounded corners=2pt, text width=7.6cm, align=center,
              minimum height=0.62cm, inner sep=3pt, fill=blue!5}]
  \node[lay] (l8) {8. Инициализация --- изграждане на приложението, регистрация};
  \node[lay, below=of l8] (l7) {7. Услуги на устройството};
  \node[lay, below=of l7] (l6) {6. Декларативен слой --- описание $\rightarrow$ обектен граф};
  \node[lay, below=of l6] (l5) {5. Контроли --- програмен интерфейс, свойства, свързване};
  \node[lay, below=of l5] (l4) {4. Разположение --- измерване и подреждане};
  \node[lay, below=of l4, fill=red!8] (l3) {3. Платформени слоеве --- свързване с native инструментариума};
  \node[lay, below=of l3, fill=green!10] (l2) {2. Обработчици --- таблици „свойство $\rightarrow$ действие“};
  \node[lay, below=of l2] (l1) {1. Договори --- виртуални изгледи, без рисуване};
  \node[lay, below=of l1] (l0) {0. Графични примитиви --- цвят, геометрия, платно};

  \draw[-Latex, thick] ([xshift=-0.45cm]l8.west) -- ([xshift=-0.45cm]l0.west)
    node[midway, rotate=90, anchor=south, yshift=2pt] {зависимостите сочат надолу};

  \draw[decorate, decoration={brace, amplitude=4pt}]
    ([xshift=0.2cm]l8.north east) -- ([xshift=0.2cm]l4.south east)
    node[midway, right=5pt, align=left] {преносимо};
  \draw[decorate, decoration={brace, amplitude=4pt}]
    ([xshift=0.2cm]l3.north east) -- ([xshift=0.2cm]l3.south east)
    node[midway, right=5pt, align=left] {платформено};
  \draw[decorate, decoration={brace, amplitude=4pt}]
    ([xshift=0.2cm]l2.north east) -- ([xshift=0.2cm]l0.south east)
    node[midway, right=5pt, align=left] {преносимо};
\end{tikzpicture}
\caption{Слоеста структура и посока на зависимостите. Само слой~3 е обвързан с native
инструментариум; тази граница прави останалите слоеве преносими и проверими без екран.}
\label{fig:layers}
\end{figure}

\subsection{Проектиране на слоя за графични примитиви}
\label{subsec:graphics_design}

Слоят съдържа стойностните типове, с които говорят всички останали слоеве: цвят, точка,
размер в дробни координати, правоъгълник, радиус на закръгление, матрица на преобразуване,
път и четки (плътна, линейна и радиална градиентна, шарка). Той е напълно преносим --- без
платформена зависимост --- затова изграждането започва от него: всеки следващ слой се
нуждае от неговите типове, той --- от нищо.

Абстракцията за рисуване (\code{canvas}) е интерфейс за команди --- линия, път, запълване,
текст --- не за конкретна повърхност. Наред с платформените реализации съществува и
\term{записваща}: не рисува, а запазва получените команди, което позволява да се провери
\emph{какво} е поискано, без екран.

Пример: превръщането на цвят от текстов запис в стойност е \code{constexpr} --- изчислено
при компилация, а в оригинала задължително при изпълнение --- така че в артефакта попада
готовата стойност, без разбор при пускане и без възможност невалиден запис да оцелее до
изпълнение.

\subsection{Договори на виртуалните изгледи}
\label{subsec:contracts}

\term{Виртуален изглед} е обект със свойства, без рисуване. Йерархията има три нива:
\code{i\_element} (най-общият елемент --- родител, идентификатор, участие в системата от
свойства, но \emph{без} размер), \code{i\_view} (добавя геометрия: желан размер,
разположение, видимост, прозрачност, отстъпи) и специализираните договори
(\code{i\_button}, \code{i\_label}, \code{i\_layout}, …), които добавят само своето.
Договорът описва \emph{какво} има контролата, никога \emph{как} изглежда --- позволява
един и същ договор да бъде удовлетворен от разнородни native контроли.

Разделянето на „елемент“ от „изглед“ не е излишно: не всеки елемент от дървото има размер.
Дефиниция на стил, елемент от менюто, колона в решетка участват в дървото, но не се
измерват. На най-общото ниво всеки от тях щеше да носи неизползвани свойства и --- по-лошо
--- да може да получи размер.

\begin{figure}[H]
\centering
\begin{tikzpicture}[font=\footnotesize, node distance=0.5cm,
  n/.style={draw, rounded corners=2pt, align=center, inner sep=4pt, minimum height=0.65cm},
  ar/.style={-Latex, thick}]
  \node[n, fill=blue!6, text width=5.4cm] (el) {\code{i\_element}\\\scriptsize родител, идентификатор, свойства --- \emph{без размер}};
  \node[n, fill=green!10, text width=5.4cm, below=of el] (vw) {\code{i\_view}\\\scriptsize размер, разположение, видимост, отстъпи};
  \node[n, text width=2.3cm, below left=0.7cm and 1.6cm of vw] (b) {\code{i\_button}};
  \node[n, text width=2.3cm, below=0.7cm of vw] (l) {\code{i\_label}};
  \node[n, text width=2.3cm, below right=0.7cm and 1.6cm of vw] (ly) {\code{i\_layout}};
  \node[n, fill=gray!8, text width=2.6cm, right=1.3cm of el] (other) {стил, елемент от меню, колона\\\scriptsize елементи БЕЗ геометрия};
  \draw[ar] (vw) -- (el); \draw[ar] (b) -- (vw); \draw[ar] (l) -- (vw); \draw[ar] (ly) -- (vw);
  \draw[ar] (other) -- (el);
\end{tikzpicture}
\caption{Йерархия на договорите: геометрията е на \code{i\_view}, не на \code{i\_element}
--- елементите вдясно нямат размер и не могат да получат такъв.}
\label{fig:contracts}
\end{figure}

\subsection{Проектиране на слоя за преобразуване}
\label{subsec:handler_design}

Обработчикът е повтарящата се единица на системата: за всяка контрола свързва виртуалния
изглед с native контролата и пренася стойностите между тях.

\subsubsection{Таблица „свойство\,$\rightarrow$\,действие“}

Преобразуването е \term{таблично}: обработчикът притежава карта от свойство към функция,
прилагаща текущата му стойност върху native контролата, и втора карта за императивните
операции. Смяна на стойност предизвиква съответното действие; свързване на нов изглед ---
пълно преминаване през таблицата.

Картите се \term{верижат} --- на контрола, на изглед, на елемент --- и разрешаването на
припокриване заслужава внимание: ключ, предефиниран в по-близка карта, надделява по
\emph{съдържание}, но се изпълнява на \emph{позицията}, на която стои в най-далечната
карта, където се среща. Редът на прилагане остава стабилен --- действие, което трябва да
предхожда друго (например установяване на контейнера), не се измества назад само защото
по-близък слой е предефинирал друго свойство. Пример: бутон, предефиниращ
\code{background}, изпълнява \emph{своята} реализация, но на позицията на
\code{background} в таблицата на изгледа --- преди \code{text}.

Съзнателни опростявания спрямо оригинала (без кеш на слетите таблици, без ключ за пакетно
прилагане) не променят \emph{кои} действия и в \emph{какъв ред}, само колко работа се
повтаря.

\subsubsection{Обобщен базов клас без виртуални шаблони}

Базата на обработчиците е шаблон, параметризиран с производния тип, типа на виртуалния
изглед и на native контролата --- вместо виртуален базов клас с шаблонни методи,
защото шаблонен метод не може да бъде виртуален: обобщената част се разрешава при
компилация, полиморфизмът --- през нешаблонен интерфейс. Незадължителните точки за
разширение се откриват чрез изискване за наличие на израз, така че производният клас
изброява само онова, което наистина предоставя. Интерфейсът на обработчика се наследява
\term{виртуално}, защото обработчик на контейнер трябва едновременно да удовлетворява
общия интерфейс и интерфейса за контейнери --- иначе би получил два подобекта на общия
интерфейс.

\subsubsection{Жизнен цикъл и собственост}

Жизненият цикъл има четири състояния --- \emph{несвързан}, \emph{свързващ се},
\emph{свързан}, \emph{пресвързващ се} --- показани на \figref{fig:handler_lifecycle};
повторно подаване на същия изглед е операция без действие. Състоянието „свързващ се“ е
нарочно \emph{видимо} за платформените действия: позволява им да пропуснат работа по
свойства със стойност и без това подразбираща се, което при първото преминаване е
повечето от тях --- без този признак всяко свързване би изпълнило пълния набор действия,
включително онези, които не променят нищо.

\begin{figure}[H]
\centering
\begin{tikzpicture}[font=\footnotesize, node distance=1.9cm,
  st/.style={draw, rounded corners=3pt, align=center, inner sep=4pt, minimum height=0.8cm,
             text width=2.1cm}, ar/.style={-Latex, thick}]
  \node[st, fill=gray!10] (d) {несвързан};
  \node[st, fill=orange!12, right=of d] (c) {свързващ се};
  \node[st, fill=green!12, right=of c] (o) {свързан};
  \node[st, fill=orange!12, below=2.0cm of o] (r) {пресвързващ се};
  \draw[ar] (d) -- node[above,font=\scriptsize,align=center] {създай\\native} (c);
  \draw[ar] (c) -- node[above,font=\scriptsize,align=center] {пълно\\преминаване} (o);
  \draw[ar] (o) -- node[right,font=\scriptsize,xshift=2pt] {нов изглед} (r);
  \draw[ar] (r.west) -| node[pos=0.32,below,font=\scriptsize] {повторно свързване} (c.south);
  \draw[ar] (o.north) -- ++(0,0.75) -| node[pos=0.25,above,font=\scriptsize] {разрушаване} (d.north);
\end{tikzpicture}
\caption{Състояния на обработчика. „Свързващ се“ е видимо за платформените действия, за да
пропуснат работа по свойства със стойност по подразбиране --- при първото преминаване
това е повечето от тях.}
\label{fig:handler_lifecycle}
\end{figure}

Собствеността е насочена: изгледът притежава обработчика, обработчикът --- native
контролата; обратните връзки са непритежаващи, следствие от липсата на автоматично
събиране на боклука (§\ref{subsec:api_design}).

Капан, установен в хода на работата: контрола с обработчик изисква вписване в \emph{два}
списъка --- регистъра на типовете и списъка, изграждан при инициализация. Вписана само в
единия, тя се компилира и стартира без грешка, но не се изобразява: липсващата регистрация
е данна, не тип, и компилаторът не може да я улови.

\subsection{Проектиране на алгоритмите за разположение}
\label{subsec:layout_design}

Разположението се изчислява в преносимия слой, не се възлага на native инструментариума
--- определящо решение: платформената контрола се третира като правоъгълник с даден
размер, не като участник в собствената си подредба. Цената е точно възпроизвеждане на
алгоритмите; ползата --- едно и също разположение навсякъде, проверимо без екран.

\subsubsection{Двуфазният модел}

Интерфейсът на един алгоритъм се състои от две операции: \emph{измерване} връща желания
размер по дадени ограничения за ширина/височина, \emph{подреждане} разполага децата по
дадени граници и връща действително заетия размер. Разделянето е необходимо, защото
размерът на родителя зависи от децата, а положението на децата --- от окончателния размер
на родителя (\figref{fig:measure_arrange}).

Ограничението „неограничено“ се предава като безкрайност --- „вземи толкова, колкото ти
трябва“ --- и се използва от подвижните области, които не ограничават децата си по
посоката на превъртане: дете, което иначе разпределя оставащото място пропорционално, тук
се свежда до собствения си желан размер. През вложени контейнери двете фази се
разпространяват рекурсивно; дълбочината на дървото умножава цената --- най-натоварената
част от преносимия слой, мястото на очакваната разлика по хипотеза H2b-2.

\begin{figure}[H]
\centering
\begin{tikzpicture}[font=\footnotesize, node distance=0.6cm,
  b/.style={draw, rounded corners=2pt, align=center, inner sep=4pt, text width=3.5cm,
            minimum height=0.7cm}, ar/.style={-Latex, thick}]
  \node[b, fill=blue!6] (m1) {родител: ограничения\\(ширина, височина)};
  \node[b, fill=blue!6, below=of m1] (m2) {дете: желан размер\\\emph{плюс} външния отстъп};
  \node[b, fill=blue!6, below=of m2] (m3) {родител: сумарен\\желан размер};
  \node[b, fill=green!10, right=2.4cm of m1] (a1) {родител: окончателни\\граници};
  \node[b, fill=green!10, below=of a1] (a2) {дете: рамка\\\emph{минус} външния отстъп};
  \node[b, fill=green!10, below=of a2] (a3) {заета област};
  \draw[ar] (m1) -- (m2); \draw[ar] (m2) -- (m3);
  \draw[ar] (a1) -- (a2); \draw[ar] (a2) -- (a3);
  \draw[ar] (m3.east) .. controls +(1.2,0) and +(-1.2,0) .. (a1.west);
  \node[above=1pt of m1, font=\bfseries] {измерване};
  \node[above=1pt of a1, font=\bfseries] {подреждане};
  \node[note, below=0.55cm of a3, xshift=-2.1cm, text width=7.4cm]
    {Отстъпът се \emph{прибавя} при измерване, \emph{изважда} се при подреждане; при нулев
     отстъп двете се компенсират --- несиметрична реализация е невидима в най-простия
     случай (§\ref{subsec:taxonomy}).};
\end{tikzpicture}
\caption{Двуфазният модел на разположението и симетрията на външния отстъп между двете фази.}
\label{fig:measure_arrange}
\end{figure}

\subsubsection{Разпределяне на пропорционалните измерения}

Разпределянето на измеренията в решетка е най-сложният алгоритъм в системата --- на
порядък по-обемен от всеки друг мениджър на разположение. Трите вида измерение се
разрешават в строг ред:

\begin{enumerate}[topsep=2pt,itemsep=1pt,leftmargin=*]
\item \textbf{Абсолютно} --- стойността е дадена; разрешава се веднага.
\item \textbf{Автоматично} --- размерът е максимумът от желаните размери на клетките в
      реда/колоната; изисква измерване на децата и затова следва абсолютните.
\item \textbf{Пропорционално} --- оставащото място след първите две се разпределя между
      тях по тегло.
\end{enumerate}

Усложнението идва от клетки, разпростиращи се върху няколко реда/колони: желаният им
размер се разпределя между обхванатите по различно правило --- пропорционалните поемат
излишъка, ако ги има сред тях, иначе го поемат автоматичните. Разпределянето се извършва
\emph{отново} при подреждане, защото окончателните граници може да се различават от
измерените.

\begin{figure}[H]
\centering
\begin{tikzpicture}[font=\footnotesize, node distance=0.55cm,
  s/.style={draw, rounded corners=2pt, align=center, inner sep=4pt, text width=3.9cm,
            minimum height=0.62cm}, ar/.style={-Latex, thick}]
  \node[s, fill=blue!6] (a) {1. абсолютни: стойността е дадена};
  \node[s, fill=blue!6, below=of a] (b) {2. автоматични: max от желаните размери на клетките};
  \node[s, fill=blue!6, below=of b] (c) {3. пропорционални: остатъкът по тегло};
  \node[s, fill=orange!12, below=of c, text width=6.6cm] (d)
     {клетки през няколко реда/колони:\\
      има ли пропорционални сред тях --- те поемат излишъка;\\ ако няма --- автоматичните};
  \node[s, fill=green!12, below=of d, text width=6.6cm] (e)
     {повторно разпределяне при подреждане\\\scriptsize (окончателните граници може да се различават)};
  \draw[ar] (a)--(b); \draw[ar] (b)--(c); \draw[ar] (c)--(d); \draw[ar] (d)--(e);
\end{tikzpicture}
\caption{Разрешаване на измеренията в решетка. Редът е задължителен: пропорционалните се
разрешават едва след абсолютните и автоматичните.}
\label{fig:grid_resolve}
\end{figure}

\subsubsection{Координатна система при подреждане}

Изборът на координати при подреждане --- абсолютни или относителни спрямо контейнера ---
изглежда като подробност и не е такава: при вложен контейнер с ненулево начало погрешният
избор води до двукратно прибавяне на отместването, проявяващо се само при влагане (клас в
систематизацията, §\ref{subsec:taxonomy}).

\subsection{Проектиране на системата от свойства}
\label{subsec:binding_design}

Системата от свойства съхранява стойностите, известява за промяна и разрешава
предимството между едновременно валидни източници на една и съща стойност.

\subsubsection{Типизиран описател вместо изтрит тип}

Описателят на едно свойство (име, стойност по подразбиране, обратни извиквания) е
\term{шаблон по типа на стойността}, за разлика от оригинала, изграден върху общия базов
тип на средата за изпълнение. Следствие: стойността по подразбиране, сравнението за
равенство и обратните извиквания са типизирани, без опаковане на стойностни типове и без
информация за типа при изпълнение.

\subsubsection{Предимство на стойностите като едно цяло число}

Предимството е представено като опакована в едно 64-битово число ключова стойност, чийто
\emph{целочислен ред съвпада с реда на предимството}. Възходящо:

\begin{figure}[H]
\centering
\begin{tikzpicture}[font=\scriptsize, node distance=0.1cm,
  lvl/.style={draw, rounded corners=2pt, align=center, inner sep=2pt,
              minimum height=0.75cm, text width=1.62cm}]
  \node[lvl, fill=gray!8]   (a) {подразбиращо се};
  \node[lvl, fill=blue!5,  right=of a] (b) {стил};
  \node[lvl, fill=blue!8,  right=of b] (c) {свързване};
  \node[lvl, fill=blue!12, right=of c] (d) {динамичен\\ресурс};
  \node[lvl, fill=green!12, right=of d] (e) {ръчно};
  \node[lvl, fill=orange!14, right=of e] (f) {тригер};
  \node[lvl, fill=orange!22, right=of f] (g) {визуално\\състояние};
  \node[lvl, fill=red!12,  right=of g] (h) {обработчик};
  \draw[-Latex, thick] ([yshift=-0.42cm]a.south west) -- ([yshift=-0.42cm]h.south east)
     node[midway, below, font=\scriptsize] {растящо предимство};
  \node[above=0.1cm of a, font=\scriptsize] {най-ниско};
  \node[above=0.1cm of h, font=\scriptsize] {най-високо};
\end{tikzpicture}
\caption{Нива на предимство при разрешаване на стойност. Подредбата им \emph{е}
целочислена. Стойността от обработчика стои най-високо, но се измества от всяко изрично
присвояване --- отразява какво е направила native контролата, не какво е поискало
приложението.}
\label{fig:precedence}
\end{figure}

Решението има две следствия: сравнението на две предимства е едно целочислено сравнение, а
стойността е изчислима при компилация, така че наименуваните константи се инициализират
статично, без неопределения ред между единиците за превод.

\subsubsection{„Незададено“ срещу „зададено на стойността по подразбиране“}

Разграничението изглежда академично, но има пряк визуален ефект: контрола с незададен цвят
на текста трябва да получи
подразбиращия се на платформата, докато контрола с изрично зададен цвят, съвпадащ с
подразбиращия се, трябва да получи точно него. Проверка на стойността вместо на наличието
ѝ обърква двата случая --- грешката, която прави текст невидим при тъмна тема (клас в
§\ref{subsec:taxonomy}).

\subsection{Проектиране на потребителския програмен интерфейс}
\label{subsec:api_design}

Липсата на автоматично събиране на боклука е най-същественото разминаване спрямо
оригинала, защото дървото от изгледи съдържа обратни връзки
родител\,$\leftrightarrow$\,дете и абонаменти за събития. Приетият модел разпределя
работата между четири роли:

\begin{center}
\begin{tabular}{@{}L{3.1cm}L{3.4cm}L{7.2cm}@{}}
\textbf{Роля} & \textbf{Средство} & \textbf{Семантика} \\
\hline
конструиране & изграждащи функции & връщат притежаващ манипулатор върху поддървото \\
притежаване & \code{view\_ref<T>} & само преместваем; точно един притежател \\
съхраняване & \code{weak\_ref<T>} & копируем, непритежаващ; безопасен достъп \\
изменение & \code{observable<T>} & реактивно; изменя се \emph{данните}, не изгледът \\
\end{tabular}
\end{center}

Решението притежаващият манипулатор да \emph{не} е копируем прави невъзможен най-често
допусканият клас грешки: улавяне на притежаващ манипулатор в обработчик, който самата
контрола притежава. Такъв цикъл не се съобщава от нито един инструмент, защото обектите са
достижими --- само един от друг; при копируем манипулатор това е валиден код, тук
е грешка при компилация. Цената в ергономичност е реална, но ограничена до достъпа през
наблюдаваща препратка с проверка за живост (измерена в §\ref{subsec:ownership}).

Където съвместното притежаване е необходимо (два елемента с независим живот,
които взаимно се изменят), съществува изричен метод, създаващ втори притежател --- изричен
и откриваем в типа: наблюдаваща препратка обявява „наблюдател“, притежаващ манипулатор ---
„притежател“.

\subsection{Формат на декларативното описание}
\label{subsec:xaml_design}

Входът е дървовидно описание на елементите с атрибути. Преобразуването минава през три
стъпки: \emph{разбор} (текст $\rightarrow$ дърво от възли), \emph{изграждане} (възли
$\rightarrow$ обекти, чрез регистър на типовете) и \emph{свързване} (изразите за свързване
стават връзки към данни). Първите две се извършват при компилация, третата --- при
изпълнение, защото зависи от данните. Регистърът на типовете и на преобразувателите (текст
$\rightarrow$ стойност) е явен и замества рефлексията: всеки тип, който може да се появи в
описанието, е вписан изрично.

\subsection{Проектиране на платформените слоеве}
\label{subsec:platform_design}

Всеки платформен слой реализира едни и същи обработчици върху различен native
инструментариум. Реализирани са пет: два върху UIKit (мобилният и настолният му вариант),
един върху native настолния инструментариум на същата операционна система, един върху
Android View системата, един върху съвременния настолен инструментариум на Windows, и
един \term{без екран}.

Слоят без екран не е спомагателен детайл --- условието, при което поведението на
преносимата част (разположение, предимство, известяване) се проверява автоматизирано, без
устройство и без изобразяване. Без него проверимостта би зависела от заснемане на
изображения.

Двата слоя върху една и съща операционна система споделят обвързаното с езика и с общите
графични услуги (цветове, шрифтове, рисуване върху платно, дати), но не и
\emph{йерархията на контролите}: различни базови класове, правила за влагане, жизнени
цикли на прозореца.

Слоят върху native настолния инструментариум е несравним попикселно с останалите ---
изгражда интерфейса от друго семейство контроли, отколкото еталонът на същата операционна
система. Изискването към него е различно (пълнота на съдържанието, съгласуваност между
двата входни пътя) и той не участва в количествените таблици (§\ref{subsec:exp_setup}):
число без смисъл е по-лошо, отколкото да липсва число.
